Validation compares the numerical trajectory with the analytical reference at
the same time nodes. The implemented error indicator is an RMS-style metric in
\texttt{src/EstimateError.m}; here, RMS stands for root mean square and measures
the typical size of the pointwise deviation between the exact and numerical
solutions.
\begin{equation}
\mathrm{RMS}
= \sqrt{\frac{1}{N-1}\sum_{i=1}^{N-1}\left(T(t_i)-\bar{T}_i\right)^2},
\label{eq:rms-error}
\end{equation}
where $T(t_i)$ denotes the exact value and $\bar{T}_i$ the numerical value.
\Cref{eq:rms-error} is reported by \texttt{RunCoolingLaw} and used
for refinement checks.

Automated tests are written with \texttt{matlab.unittest} in
\texttt{tests/TestCoolingLaw.m}. The suite checks:
\begin{itemize}
\item exact-solution consistency at representative points,
\item numerical output dimensions and boundary consistency,
\item monotonic cooling behavior for default parameters,
\item convergence trends under grid refinement,
\item regression snapshots for reproducible default behavior.
\end{itemize}

Together, these checks provide both mathematical confidence and practical
regression protection during refactoring. The visual convergence evidence in
\Cref{fig:trajectory-convergence} and \Cref{fig:endpoint-convergence}
complements the automated checks.
